%!TEX root =  tfg.tex
\chapter{Planificación}

\begin{quotation}[Novelist]{Ernest Hemingway (1899--1961)}
The good parts of a book may be only something a writer is lucky enough to overhear or it may be the wreck of his whole damn life -- and one is as good as the other.
\end{quotation}

\begin{abstract}
Este capítulo describe la estrategia de gestión y el cronograma seguidos 
para el desarrollo de \textit{Via Inferni}. Se detalla la división del proyecto en 
cinco fases lógicas, desde la concepción inicial hasta el cierre, justificando el 
esfuerzo temporal y la carga de trabajo necesaria para cumplir con los estándares de un 
Trabajo de Fin de Grado en Ingeniería Informática. El objetivo es ofrecer una visión clara 
de cómo se han administrado los recursos y el tiempo para transformar los requisitos
en un producto funcional.
\end{abstract}

\section{Resumen temporal del proyecto}

\begin{table*}[htb]
	\centering
	\begin{coolTable}{ll}{2}
{Resumen del proyecto}
	\textbf{Fecha de inicio}&10/10/2014\\
	\textbf{Fecha de fin}&10/10/2014\\
	\textbf{Periodicidad de las revisiones}&3 semanas\\
	\textbf{Carga de trabajo semanal}&12 horas\\
	\textbf{Horas totales previstas}&225 horas\\ % entre 25-30 horas por crédito
	\textbf{Horas finales}&234 horas\\
	\end{coolTable}
	\caption{Tabla resumen de tiempos y planificación}
\end{table*}

\section{Planificación inicial}

\textbf{ESTE CAPÍTULO DEBE ESCRIBIRSE AL COMIENZO DEL PROYECTO}

En esta sección se presenta el desglose inicialmente planificado de las fases en las que se ha dividido el proyecto,
así como una justificación de cada una de ellas. Se han definido cinco fases principales:

\begin{table*}[htb]
	\centering
	\begin{coolTable}{ll}{2}
{Resumen de fases del proyecto}
	\textbf{Fase 1: Planificación y Diseño}&09/02/2026 a 15/02/2026\\
	\textbf{Fase 2: Implementación del Núcleo}&16/02/2026 a 15/03/2026\\
	\textbf{Fase 3: Implementación de Funcionalidades}&16/03/2026 a 19/04/2026\\
	\textbf{Fase 4: Pruebas, Depuración y Mejoras}&20/04/2026 a 10/05/2026\\
	\textbf{Fase 5: Cierre del Proyecto}&11/05/2026 a 25/05/2026\\
	\end{coolTable}
	\caption{Planificación temporal de fases del proyecto}
\end{table*}

El trabajo se realizará del 9 de febrero al 25 de mayo, con una duración total de aproximadamente 
tres meses y medio, y un total de 99 días. Se estiman unas 3 o 4 horas de trabajo diario,
lo que se traduce en un total aproximado de 300 horas de trabajo por miembro. 

Cada integrante adaptará las horas de trabajo diarias a su disponibilidad, pero se mantendrá 
una comunicación constante con reuniones recurrentes para asegurar la integración. Al menos
una vez al final de cada semana se realizará una reunión de seguimiento para revisar el progreso y
planificar las tareas de la semana siguiente.

\subsection{Fase 1: Planificación y Diseño}

Esta fase consistirá en realizar un análisis detallado de los requisitos del proyecto, definir 
los objetivos específicos y elaborar un plan de trabajo detallado. Además, habrá que diseñar 
la arquitectura del sistema y seleccionar las tecnologías y herramientas a utilizar.

Se ha estimado una duración de una semana para esta fase, considerando el tiempo necesario para 
la investigación, la elaboración de documentos de diseño y la planificación detallada del proyecto.
Además, todo el trabajo previo hace posible esta aparentemente corta planificación, ya que se han realizado
tareas de investigación y recopilación de información durante el periodo previo al inicio oficial del proyecto.

\subsection{Fase 2: Implementación del Núcleo}

Durante esta fase se llevará a cabo la implementación de los componentes fundamentales del sistema,
lo que sería el core del juego. Esto es el sistema de generación procedural de niveles, la gestión 
de personajes, enemigos y objetos, y en general la implementación de la mecánica de juego básica.

El objetivo de esta fase es establecer una base sólida sobre la cual se puedan construir las funcionalidades
adicionales en fases posteriores, como podrán ser una mayor variedad de niveles, nuevos objetos y diferentes
patrones de comportamiento de enemigos.

La implementación del núcleo se llevará a cabo de forma conjunta por los dos integrantes del equipo,
de forma que se facilite el trabajo y la colaboración. Las tareas durante la siguiente fase podrán
ser más individualizadas, pero durante esta fase se ha considerado importante que ambos integrantes 
estén involucrados en la implementación del núcleo para asegurar una comprensión compartida del sistema. 

Se ha estimado una duración de cuatro semanas para esta fase, considerando la complejidad de los
componentes a implementar y la falta de experiencia previa en el desarrollo de este tipo de sistemas.
La idea es que ambos integrantes del equipo puedan colaborar y compartir el conocimiento, facilitando
el trabajo de la siguiente fase.

\subsection{Fase 3: Implementación de Funcionalidades}

En esta fase se implementarán las funcionalidades adicionales del sistema, como la generación de niveles 
más variados, la incorporación de nuevos objetos y enemigos, y la mejora de la mecánica de juego básica 
establecida en la fase anterior.

Como ya se ha mencionado, las tareas en esta fase podrán ser más individualizadas. Esto permitirá
que, en caso de que surja algún tipo de dificultad que provoque que un miembro tenga que dejar
de trabajar, o una eventual separación del equipo, cada integrante podrá mostrar y justificar
su trabajo individual. De cualquier forma, y teniendo en cuenta que lo anteriormente mencionado
se trata de una situación excepcional, se mantendrá una comunicación constante para asegurar 
la coherencia del sistema y la integración.

Se ha estimado una duración de cinco semanas para esta fase, teniendo en cuenta que la cantidad de funcionalidades
adicionales que se pueden desarrollar es potencialmente muy amplia y que se pueden presentar desafíos técnicos
durante la implementación. El alcance final de esta fase se irá adaptando en función del progreso realizado y
de las dificultades encontradas, pero se ha considerado que cinco semanas es un plazo razonable para implementar 
una cantidad significativa de funcionalidades adicionales. Cabe destacar que una de estas semanas, la del 
29 de marzo al 29 de abril, coincide con la Semana Santa, festividad que puede afectar el ritmo de 
trabajo del equipo.

\subsection{Fase 4: Pruebas, Depuración y Mejoras}

En esta fase se llevarán a cabo pruebas exhaustivas del sistema para identificar y corregir errores, 
así como para mejorar el rendimiento y la experiencia de usuario. Se realizarán pruebas unitarias, 
pruebas de integración y pruebas de aceptación para asegurar que el sistema cumple con los requisitos 
establecidos.

Se ha estimado una duración de tres semanas para esta fase, considerando el tiempo necesario para realizar
las tareas de prueba, depuración y mejora. Además, esta fase es crucial para garantizar la calidad del 
producto final, por lo que es tiempo adecuado para llevar a cabo estas actividades de manera exhaustiva.
También hay que destacar que del 21 al 26 de abril se celebra la Feria de Sevilla, festividad que puede afectar 
el ritmo de trabajo del equipo.

\subsection{Fase 5: Cierre del Proyecto}

En esta fase se realizará la documentación final del proyecto, incluyendo el manual de usuario y el 
informe final. Además, se preparará la presentación final del proyecto y se llevará a cabo una revisión 
general para asegurar que todo el trabajo realizado esté correctamente documentado y presentado.

Se ha estimado una duración de dos semanas para esta fase, considerando que la memoria se encontrará en
un estado avanzado al trabajar sobre ella a lo largo de todo el proyecto, y que la preparación de 
la presentación final requerirá tiempo para organizar el contenido y ensayar la presentación.

\section{Informe de tiempos del proyecto}

Lo mismo que el anterior pero con datos reales. Ver Tabla \ref{tab:InformeTiempos}.

\begin{table*}[htb]
	\centering
	\begin{coolTable}{ll}{2}
{Resumen de iteraciones}
	\textbf{Iteración 1}&10/10/14 a 21/10/14\\
	\textbf{Iteración 2}&21/10/14 a 15/11/14\\
	\textbf{...}&dd/mm/aa a dd/mm/aa\\
	\end{coolTable}
	\caption{Planificación temporal de iteraciones\label{tab:InformeTiempos}}
\end{table*}

Justificar los retrasos de forma detallada aquí para cada una de las iteraciones. Explicar las razones.